\section{Hai phép điều chuẩn, và một dự đoán sai}
\label{sec:ch08-dieu-chuan}

\index{điều chuẩn!mục 5.4 nói ba, in ra hai}%
Mục 5.4 của bài mở bằng \enquote{We employ three types of regularization during
training:} rồi in ra hai tiêu đề đậm, \code{Residual Dropout} và
\code{Label Smoothing}. Tôi tìm cái thứ ba ở cả hai bản và không có. Cách đọc
duy nhất làm chữ \enquote{three} đúng là đếm đoạn dropout thành hai, vì nó nêu
hai chỗ đặt khác nhau:

\begin{quote}
\enquote{We apply dropout to the output of each sub-layer, before it is added
to the sub-layer input and normalized. In addition, we apply dropout to the sums
of the embeddings and the positional encodings in both the encoder and decoder
stacks. For the base model, we use a rate of \(P_{drop} = 0.1\).}
\end{quote}

Chuyện đếm ấy nhỏ. Chuyện đáng nói là câu ngay sau đó, về phép thứ hai:

\begin{quote}
\enquote{During training, we employed label smoothing of value
\(\epsilon_{ls} = 0.1\). This hurts perplexity, as the model learns to be more
unsure, but improves accuracy and BLEU score.}
\end{quote}

\index{label smoothing!bài tự nói nó làm xấu perplexity}%
Một bài báo nói thẳng rằng thứ họ dùng làm xấu một độ đo của chính họ. Đó là câu
làm tôi muốn đo mục này, vì nó là một phát biểu kiểm được, và bảng 3 của chính
bài kiểm nó.

\subsection*{Bảng 3 hàng (D), và hai độ đo đi ngược nhau}

\index{điều chuẩn!hàng (D) của bảng 3}%
Hàng (D) đổi hai tham số ấy quanh cấu hình base, và đo trên tập phát triển
Anh-Đức. Cấu hình base là \(P_{drop} = 0.1\), \(\epsilon_{ls} = 0.1\), với
perplexity 4.92 và BLEU 25.8:

\begin{minted}{text}
                PPL (dev)   BLEU (dev)
base            4.92        25.8
P_drop  = 0.0   5.77        24.6
P_drop  = 0.2   4.95        25.5
eps_ls  = 0.0   4.67        25.3
eps_ls  = 0.2   5.47        25.7
\end{minted}

Đọc từng nửa.

\textbf{Dropout.} Tắt hẳn thì perplexity đi từ 4.92 lên 5.77 và BLEU từ 25.8
xuống 24.6. Hai độ đo cùng xấu đi, và 1.2 BLEU là một khoảng lớn theo tiêu chuẩn
của chính bảng ấy, nơi cả bảng chỉ trải từ 23.7 tới 26.4. Đẩy lên 0.2 thì cũng
xấu hơn base nhưng ít. Nên 0.1 là một cực trị đã chỉnh, không phải một con số
mặc định.

\textbf{Label smoothing.} Tắt hẳn thì perplexity đi từ 4.92 \emph{xuống} 4.67,
tức \emph{tốt hơn}, trong khi BLEU đi từ 25.8 xuống 25.3, tức xấu hơn. Hai cột
đi ngược chiều nhau, đúng bằng câu ở mục 5.4, và đo được trong bảng của chính
bài.

Đó là chỗ chữ \enquote{mẹo} hỏng hẳn. Một mẹo là thứ làm mọi độ đo tốt lên một
chút. Cái này làm một độ đo xấu đi để đổi lấy một độ đo khác, và nhóm tác giả
biết điều đó, viết ra, rồi vẫn giữ. Đó là một quyết định về việc độ đo nào là độ
đo thật.

\begin{bridge}{Szegedy và cộng sự 2016: label smoothing đến từ đâu}
\index{label smoothing!Szegedy 2016}%
Bài này về kiến trúc Inception cho ảnh, không liên quan gì tới
chuỗi~\autocite{szegedy2016rethinking}. Mục 7 đưa ra phép làm mượt nhãn: thay
phân phối đích một điểm \(\delta_{k,y}\) bằng

\begin{equation*}
  q'(k \mid x) = (1 - \epsilon)\,\delta_{k,y} + \epsilon\, u(k),
\end{equation*}

với \(u(k)\) là phân phối đều, tức \(\epsilon/K\) rải đều trên \(K\) lớp. Bài
viết \enquote{In our ImageNet experiments with \(K = 1000\) classes, we used
\(u(k) = 1/1000\) and \(\epsilon = 0.1\)}. Đúng \(\epsilon\) mà Vaswani và cộng
sự lấy lại một năm rưỡi sau, cho một bài toán khác hẳn.

Hiệu quả bài báo cáo: \enquote{a consistent improvement of about 0.2\% absolute
both for top-1 error and the top-5 error}. Chữ \enquote{about} nên trích kèm,
vì bảng 3 của chính bài cho 0.3 điểm phần trăm ở top-1 và 0.2 ở top-5, không
phải 0.2 cho cả hai.

Tôi đọc bài này ở riêng mục 7.
\end{bridge}

\subsection*{Dự đoán trước khi chạy}

\index{điều chuẩn!dự đoán trước khi chạy}%
Corpus của sách này không phải WMT'14, và tôi viết dự đoán ra trước khi chạy để
sau đó không viết lại nó cho khớp kết quả.

Dropout là phép phòng overfit, và bài lấy nó từ Srivastava và cộng sự
2014~\autocite{srivastava2014dropout} chứ không phải từ Hinton và cộng sự 2012,
điều đáng nói vì chỗ trích lại hay đổi hai bài cho nhau.
Mục~\ref{sec:ch07-corpus} đã đo rằng kiến trúc
này \emph{không} overfit trên corpus này: chấm mô hình \(\dmodel = 32\) của mục
ấy trên 300 cặp lấy từ chính tập huấn luyện cho 0.5967, không cao hơn 0.6533 nó
được trên tập test. Đó là một mô hình rộng hơn mô hình bảng dưới đây huấn luyện,
nên nó là bằng chứng về kiểu chứ không phải về đúng cấu hình này. Một phép phòng
một chuyện không xảy ra thì chỉ có thể tốn. Dự đoán: dropout làm hỏng.

Label smoothing thì tôi dự đoán cùng chiều, và lý do tôi đưa ra lúc ấy là văn
phạm sinh ra corpus này hữu hạn nên đích thật sự là một điểm, không phải một
phân phối. Dự đoán: cũng làm hỏng.

Bảng chạy 2x2 chứ không quét từng cái một, vì bài chưa bao giờ chạy riêng một
trong hai và tôi muốn thấy chúng có tương tác không. Ba seed mỗi ô, vì lý do
mục~\ref{sec:ch08-warmup} vừa trả giá.

\begin{measured}{dropout và label smoothing trên corpus của sách, 2x2, ba seed}
\begin{minted}[fontsize=\footnotesize]{text}
P_drop  eps_ls  loss    floor   excess  exact   min     max     long
0.0     0.0     0.1449  0.0000  0.1449  0.5333  0.4700  0.5667  0.1798
0.1     0.0     0.3644  0.0000  0.3644  0.3522  0.2867  0.4233  0.0219
0.0     0.1     0.7910  0.6678  0.1233  0.5700  0.5267  0.6067  0.1886
0.1     0.1     1.0050  0.6678  0.3372  0.3822  0.3467  0.4267  0.0175
\end{minted}

2 tầng, 14 epoch, \(\dmodel = 24\), seed 0, 1, 2. \tn{Từ điển}{vocabulary} đích
37 lớp.

Đo bằng \code{experiments/ch08\_recipe.py} tại tag \code{ch08}.
\end{measured}

\subsection*{Cột \code{floor}, và vì sao cột \code{loss} một mình nói dối}

\index{label smoothing!lower bound của mất mát}%
Trước khi đọc kết quả, phải xử lý một chuyện về thang đo, và nó là chỗ dễ kết
luận ngược nhất trong cả chương.

Dưới label smoothing, đích không còn là một điểm, nên một mô hình \emph{hoàn
hảo} vẫn phải trả mất mát. Giá trị nhỏ nhất của entropy chéo với đích đã làm
mượt là entropy của chính đích ấy:

\begin{equation}
  H(q') = -\left(1 - \epsilon + \tfrac{\epsilon}{K}\right)
            \log\!\left(1 - \epsilon + \tfrac{\epsilon}{K}\right)
          - (K-1)\,\tfrac{\epsilon}{K} \log \tfrac{\epsilon}{K}.
  \label{eq:ch08-chan-duoi}
\end{equation}

Với \(K = 37\) và \(\epsilon = 0.1\), \eqref{eq:ch08-chan-duoi} cho 0.6678, và
đó là cột \code{floor}. Cột \code{excess} là mất mát trừ lower bound, tức phần
mô hình thật sự còn thiếu.

Đọc cột \code{loss} thô thì label smoothing làm mất mát tệ đi 5.5 lần, từ 0.1449
lên 0.7910. Đọc cột \code{excess} thì nó \emph{tốt lên}, từ 0.1449 xuống 0.1233.
Cùng một phép đo, hai kết luận ngược nhau, và cái đúng là cái đã trừ lower
bound. Không có cột ấy thì bảng này chỉ nói được rằng smoothing làm một con số
to lên, mà nó làm thế theo định nghĩa.

\subsection*{Một dự đoán đúng, một dự đoán sai}

\index{điều chuẩn!dropout hỏng ngoài nhiễu}%
\textbf{Dropout làm hỏng, và vượt nhiễu.} 0.5333 xuống 0.3522 khi không có
smoothing, 0.5700 xuống 0.3822 khi có. Quan trọng hơn hai con số ấy: khoảng
\code{min}-\code{max} của hai bên không giao nhau, 0.4700 tới 0.5667 so với
0.2867 tới 0.4233. Đây là hiệu ứng thật chứ không phải seed. Câu dài chịu gần
hết, 0.1798 xuống 0.0219.

Dự đoán đúng, và lý do đứng vững: bỏ ngẫu nhiên một phần tín hiệu ở một mô hình
đang thiếu tín hiệu thì chỉ tốn. Chỗ này \emph{không} mâu thuẫn với hàng (D) của
bài. Bài đo trên 4.5 triệu cặp câu thật, nơi overfit là chuyện có thật; sách này
đo trên 6000 câu sinh từ một văn phạm hữu hạn, nơi nó không phải. Hai bảng trả
lời hai câu hỏi khác nhau, và cái chung của chúng là dropout không phải một mẹo
mà là một phép đánh đổi có điều kiện, với điều kiện là mô hình đang overfit.

\index{label smoothing!dự đoán của chương sai}%
\textbf{Label smoothing thì không làm hỏng, và ở đây tôi dự đoán sai.} Khớp đúng
cả câu đi từ 0.5333 lên 0.5700 và từ 0.3522 lên 0.3822, tức tốt lên ở cả hai
điều kiện dropout. Khoảng \code{min}-\code{max} có giao nhau nên độ lớn không
kết luận được, nhưng chiều thì nhất quán, và cột \code{excess} cũng đi cùng
chiều ấy, 0.1449 xuống 0.1233.

Chỗ tôi sai nằm ở lập luận chứ không ở số. Tôi coi label smoothing là một họ với
dropout, tức một phép phòng overfit, nên đã suy ra cùng một kết luận từ cùng một
tiền đề. Nó không phải. Dropout bỏ tín hiệu đi trong lúc học; label smoothing
giữ nguyên tín hiệu và chỉ chặn mô hình đừng đẩy xác suất của một token lên sát
1. Corpus không có gì để overfit vẫn có thể có mô hình quá tự tin, và hai chuyện
ấy độc lập với nhau.

Chiều của kết quả cũng khớp với bài theo một nghĩa hẹp nhưng thật. Bảng 3 hàng
(D) cho hai độ đo tách nhau, perplexity và BLEU. Ở đây khớp đúng cả câu là độ đo
họ BLEU, và \code{excess} là độ đo họ perplexity, và cả hai cùng nói smoothing
có lợi nhẹ. Không cùng kết quả với bài, vì bài thấy perplexity xấu đi, nhưng
cùng một chuyện: hai độ đo của cùng một mô hình không buộc phải đi cùng chiều,
và chọn nhìn cái nào là một quyết định chứ không phải một chi tiết.

\index{điều chuẩn!tương tác giữa hai phép}%
Còn tương tác thì gần như không có, và đó cũng là một kết quả. Dropout lấy đi
0.1811 khi không có smoothing và 0.1878 khi có; smoothing thêm 0.0367 khi không
có dropout và 0.0300 khi có. Hai hiệu ứng cộng vào nhau chứ không nhân lên,
trong phạm vi bảng này phân biệt được. Bài chưa bao giờ chạy riêng một trong hai
nên không có gì để đối chiếu.
